% Irodov, Arihant PDF pp. 36--38, printed pp. 27--29.
% Momentum conservation, ballistic pendulum, elastic/inelastic collisions.
\exproblem{161}{27}
质量为 $M$ 的火炮从静止开始沿倾角为 $\alpha$ 的
光滑斜面自由滑下，沿斜面移动距离 $\ell$ 后射出一枚
水平方向动量大小为 $p$ 的炮弹（相对于地面）。
因反冲，火炮在发射后停止。
炮弹质量相对火炮可忽略。求发射持续时间 $\tau$。
\exsolution
设火炮沿斜面向下的速度刚开始发射时为 $u$，
由发射前机械能守恒
\[
\frac12Mu^2=Mg\ell\sin\alpha
\quad\Longrightarrow\quad
u=\sqrt{2g\ell\sin\alpha}.
\]
取沿斜面向下为正。
在发射持续时间 $\tau$ 内，
重力沿斜面的冲量为 $Mg\sin\alpha\,\tau$。
为使火炮停下，炮弹水平方向发射须在斜面向下方向
具有正投影 $p\cos\alpha$，
则火炮受到反向的沿斜面冲量 $-p\cos\alpha$。
忽略炮弹相对火炮质量造成的质量修正，沿斜面动量定理：
\[
0-Mu=Mg\sin\alpha\,\tau-p\cos\alpha .
\]
所以
\[
\boxed{\tau=
\frac{p\cos\alpha-M\sqrt{2g\ell\sin\alpha}}
 {Mg\sin\alpha}}.
\]
\excheck 物理上要求 $\tau\ge0$，即
$p\cos\alpha\ge M\sqrt{2g\ell\sin\alpha}$；
若不满足，不可能以给定脉冲停止炮身。
结果在理想脉冲模型下成立，未考虑实际炮身运动时滑轮或轮轨阻力。

\exproblem{162}{27--28}
质量为 $m$ 的子弹水平射入悬挂于两根等长细绳上的
质量 $M$ 的物体中并嵌入（图1.42），随后
悬挂系统向上摆至细绳偏离竖直 $\theta$。
设 $m\ll M$、绳长为 $\ell$。
求：（a）碰前子弹速率；
（b）子弹初始动能中转化为热的比例。
\exsolution
碰撞时间很短，沿水平方向动量守恒。
碰撞后子弹与物体组成整体，瞬时速度为 $V$：
\[
mv=(M+m)V.
\]
碰撞结束到摆至最高点，机械能守恒：
\[
\frac12(M+m)V^2=(M+m)g\ell(1-\cos\theta),
\]
故
\[
V=\sqrt{2g\ell(1-\cos\theta)}
=2\sqrt{g\ell}\sin\frac\theta2.
\]
因此
\[
\boxed{v=\frac{2(M+m)}m\sqrt{g\ell}
\sin\frac\theta2
\simeq\frac{2M}{m}\sqrt{g\ell}\sin\frac\theta2}.
\]
碰撞中的机械能损失比例为
\[
\eta_{\rm heat}=
1-\frac{\tfrac12(M+m)V^2}{\tfrac12mv^2}
=1-\frac{m}{M+m}
=\boxed{\frac{M}{M+m}\simeq1-\frac mM}.
\]
\excheck 短时碰撞\emph{只守恒水平动量}，不守恒动能；
碰后摆起阶段才使用机械能守恒。

\exproblem{163}{28}
质量为 $M$ 的可水平滑动大物块置于光滑水平面，
其上放一质量为 $m$ 的小圆片，
圆片获得初始水平速率 $v$，并沿大物块的
光滑弧形表面上爬（图1.43）。
求圆片能够达到的最大高度（相对初始高度）。
\exsolution
水平外力为零，因此系统水平动量
$P_x=mv$ 守恒。初始总动能为
$K_0=\tfrac12mv^2$。
在圆片达到最大可能高度时，系统动能不能小于
同一水平动量对应的质心整体平移动能
$K_{\rm CM}=P_x^2/[2(M+m)]$。
若在上升最高位置的相对速度瞬时为零，
该下界可以达到，因而
\[
mgh_{\max}=
\frac12mv^2-\frac{m^2v^2}{2(M+m)}
=\frac{mMv^2}{2(M+m)}.
\]
于是
\[
\boxed{h_{\max}=\frac{Mv^2}{2g(M+m)}}.
\]
\excheck 这是原书理想光滑导轨装置可达到的极限高度；
$M\to\infty$ 时回到固定斜面上的 $v^2/(2g)$，
$m\gg M$ 时可达到的高度趋小。

\exproblem{164}{28}
质量为 $m$ 的小圆片从高度 $h$ 的光滑小山
由静止滑下，到达底部后进入质量为 $M$、
初始静止的水平木板（地面光滑），
圆片与木板之间存在摩擦，最终两者共同运动（图1.44）。
求：（1）此过程中摩擦力所做总功；
（2）这一结果是否依赖惯性参考系的选择。
\exsolution
圆片到山脚时速率
$u=\sqrt{2gh}$，木板静止。
进入木板后系统水平方向没有外力，
因此最终共同速度由动量守恒给出
\[
V=\frac{mu}{m+M}.
\]
摩擦耗散的总机械能等于最终与初始动能之差：
\[
W_f=\frac12(M+m)V^2-\frac12mu^2
=-\frac12\frac{mM}{m+M}u^2.
\]
用 $u^2=2gh$，得到
\[
\boxed{W_f=-\frac{mM}{m+M}gh=-\mu gh},
\qquad \mu=\frac{mM}{m+M}.
\]
（2）两物体的摩擦力大小相等、方向相反，故它们做功的
总和可以写为
$W_f=-\int f\,\mathrm d(x_{\rm disc}-x_{\rm plank})$。
两物体的相对位移在所有伽利略惯性系中相同，
所以 $\boxed{W_f\text{ 不依赖惯性系选择}}$。
\excheck 这里不变的是\emph{摩擦力的合计总功}，
并非摩擦力分别对两物体所做的单独功。

\exproblem{165}{28}
一颗石子从地面上方高度 $h$ 的地方相对地面静止释放，
忽略空气阻力。试改用一个相对地面以恒速
$v_0$ 向下运动的惯性参考系，重新得到
石子撞地速率 $v_0=\sqrt{2gh}$。
\exsolution
设所选参考系向下匀速平移，速率等于石子
\emph{最终}撞地时相对地面的速率 $v_0$。
在该系中，石子起始速度竖直向上，大小为 $v_0$；
由于重力向下，石子作加速度 $-g$ 的运动，
恰在撞地时相对该系速度为零。
这段时间 $T=v_0/g$，石子在移动系中向上位移
\[
\Delta y' =v_0T-\frac12gT^2=\frac{v_0^2}{2g}.
\]
地面在该移动系中向上走过距离
$v_0T=v_0^2/g$。
初始石子与地面相隔 $h$，撞地时位置重合，
故地面比石子多向上走过的距离恰为 $h$：
\[
h=\frac{v_0^2}{g}-\frac{v_0^2}{2g}
=\frac{v_0^2}{2g}
\quad\Longrightarrow\quad
\boxed{v_0=\sqrt{2gh}}.
\]
\excheck 这个参考系虽向下运动，却仍是\emph{惯性系}；
不能把地面在该系中的运动误当成零。

\exproblem{166}{28}
一质量为 $m_1=1.0\,\mathrm g$ 的质点
初速度 $\bm v_1=3.0\bm i-2.0\bm j\,(\mathrm{m/s})$，
与质量为 $m_2=2.0\,\mathrm g$、
初速度 $\bm v_2=4.0\bm j-6.0\bm k\,(\mathrm{m/s})$
的质点发生完全非弹性碰撞并结合。
求结合体的速度矢量及速率。
\exsolution
忽略外界冲量，总动量守恒：
\[
(m_1+m_2)\bm v=m_1\bm v_1+m_2\bm v_2.
\]
按各坐标分量计算
\[
\bm v=
\frac{1(3,-2,0)+2(0,4,-6)}{1+2}
=(1,2,-4)\,\mathrm{m/s}.
\]
因此
\[
\boxed{\bm v=(1.0\bm i+2.0\bm j-4.0\bm k)\,\mathrm{m/s}},
\qquad
\boxed{|\bm v|=\sqrt{1+4+16}=\sqrt{21}
\simeq4.58\,\mathrm{m/s}}.
\]
\excheck 分量计算使用相同质量单位 g，
故分子和分母的单位约去；碰撞后两球必有同一速度。

\exproblem{167}{28}
由质量 $m_1,m_2$ 的两球构成的孤立系统
发生完全非弹性碰撞，碰前速度分别为
$\bm v_1,\bm v_2$。求系统总动能增量。
\exsolution
碰后结合体速度由动量守恒确定：
\[
\bm V=\frac{m_1\bm v_1+m_2\bm v_2}{m_1+m_2}.
\]
动能变化为
\[
\Delta K=\frac12(m_1+m_2)V^2-
\frac12m_1v_1^2-\frac12m_2v_2^2.
\]
展开平方并合并同类项，得到
\[
\boxed{\Delta K=
-\frac12\mu|\bm v_1-\bm v_2|^2},
\qquad
\mu=\frac{m_1m_2}{m_1+m_2}.
\]
\excheck 动能变化非正，仅当两球碰前速度矢量完全相同
时为零。动量守恒不意味着动能守恒。

\exproblem{168}{28}
质量 $m_1$ 的粒子与静止质量 $m_2$ 的粒子
发生完全弹性碰撞。求入射粒子损失的初始动能比例：
（a）碰后入射粒子速度与原方向垂直；
（b）正碰。
\exsolution
设入射初速为 $v$，碰后入射速率为 $u$，
目标粒子碰后速度由动量守恒
$\bm w=(m_1/m_2)(\bm v-\bm u)$ 确定。
（a）若 $\bm v\cdot\bm u=0$，则
$w^2=(m_1/m_2)^2(v^2+u^2)$。
代入弹性碰撞动能守恒并整理：
\[
m_1v^2=m_1u^2+
\frac{m_1^2}{m_2}(v^2+u^2)
\Longrightarrow
\frac{u^2}{v^2}=\frac{m_2-m_1}{m_1+m_2}.
\]
故
\[
\boxed{\eta_a=1-\frac{u^2}{v^2}
=\frac{2m_1}{m_1+m_2}}.
\]
（b）一维正碰后入射粒子速度
$u=(m_1-m_2)v/(m_1+m_2)$，所以
\[
\boxed{\eta_b=
1-\left(\frac{m_1-m_2}{m_1+m_2}\right)^2
=\frac{4m_1m_2}{(m_1+m_2)^2}}.
\]
\excheck （a）要求 $m_2>m_1$ 才有非零速率的严格
$90^\circ$ 散射；$m_2=m_1$ 的极限入射粒子停下，
其“速度方向”不再有定义。不能让能量损失比例超过 $1$。

\exproblem{169}{28--29}
粒子1与静止粒子2发生完全弹性碰撞。求质量比 $m_1/m_2$：
（a）正碰后两粒子反向飞开且速率相等；
（b）碰后两粒子的飞行方向相对入射方向对称，
两方向夹角为 $\theta=60^\circ$。
\exsolution
（a）设碰前粒子1速度 $v$，
碰后两者速度为 $-u,+u$，其中 $u>0$。
动量与动能守恒分别给出
\[
m_1v=(m_2-m_1)u,\qquad
m_1v^2=(m_1+m_2)u^2.
\]
消去 $u,v$，得到
$(m_2-m_1)^2=m_1(m_1+m_2)$，
物理允许解
$\boxed{m_1/m_2=1/3}$。
（b）两射出方向分别相对初始方向倾斜
$+\theta/2$ 与 $-\theta/2$。
由横向动量守恒，
$m_1u_1=m_2u_2$；
由纵向动量守恒，
$m_1v=(m_1u_1+m_2u_2)\cos(\theta/2)
=2m_1u_1\cos(\theta/2)$。
因此 $u_1=v/[2\cos(\theta/2)]$。
动能守恒
$m_1v^2=m_1u_1^2+m_2u_2^2
=m_1u_1^2(1+m_1/m_2)$，
故
\[
\boxed{\frac{m_1}{m_2}=4\cos^2(\theta/2)-1
=1+2\cos\theta=2.0}.
\]
\excheck （b）要求两侧角度对称而非两粒子碰后速率相等；
横向动量守恒决定了两速率的质量加权关系。

\exproblem{170}{29}
质量相等的两光滑小球发生弹性碰撞，
其中一个小球以速率 $v$ 平动，另一个原先静止。
碰撞瞬间球心连线与入射速度成
$\alpha=45^\circ$ 角。
求两球压缩达到最大形变瞬间，
初始动能中暂时转为弹性势能的比例。
\exsolution
把入射速度分解为沿碰撞法线（球心连线）
的分量 $v_n=v\cos\alpha$，
及切向分量 $v_\tau=v\sin\alpha$。
光滑球无切向冲量，因此只在法向发生压缩。
在最大压缩瞬间，两球法向速度相等；
由法向动量守恒，两球共同法向速度为 $v_n/2$。
初始法向动能为 $\tfrac12mv_n^2$，
最大压缩时法向平动动能为
$2(\tfrac12 m(v_n/2)^2)=\tfrac14mv_n^2$，
故暂存弹性势能为
\[
U_{\rm elastic,max}=\frac14mv_n^2
=\frac14mv^2\cos^2\alpha .
\]
与初始总动能 $\tfrac12 mv^2$ 相比：
\[
\boxed{\eta=
\frac{U_{\rm elastic,max}}{\tfrac12 mv^2}
=\frac12\cos^2\alpha
=\frac14=25\%}.
\]
\excheck 完全弹性碰撞后的最终势能又恢复为零；
这里所问的是\emph{最大压缩瞬间}，不是碰后耗散热的比例。
